\documentclass{article} 
\usepackage{graphicx}
\usepackage[francais]{babel} 
\usepackage[utf8]{inputenc}
\usepackage[T1]{fontenc} 
\usepackage{amsmath} 
\usepackage{amsfonts}
\usepackage{verbatim} 
\usepackage{float} 
\usepackage{hyperref}
\usepackage{scrtime}
\usepackage[scale=2]{ccicons}
\usepackage{tabularx}
\usepackage{Sweave}
\begin{document}
%%% Police de caractère pour la page titre
\renewcommand{\sfdefault}{hls}

%%% Marge de gauche (1/3 de la page)
\newlength{\gauche}
\addtolength{\gauche}{72mm}
\addtolength{\gauche}{-\spinemargin}

%%% Épaisseur de la bande sur la page titre
\newlength{\ruleheight}
\setlength{\ruleheight}{7.75mm}

%%% Définition de la bande
\definecolor{gray}{gray}{0.4}
\textblockorigin{0mm}{279mm}
\newcommand{\banderecto}{%
  \begin{textblock*}{71.5mm}[0,1](0mm,-46.5mm)
    \textblockcolor{black} \rule{0mm}{\ruleheight}
  \end{textblock*}
  \begin{textblock*}{144mm}[0,1](72mm,-46.5mm)
    \textblockcolor{gray} \rule{0mm}{\ruleheight}
  \end{textblock*}}
\newcommand{\bandeverso}{%
  \begin{textblock*}{144mm}[0,1](0mm,-46.5mm)
    \textblockcolor{gray} \rule{0mm}{\ruleheight}
  \end{textblock*}
  \begin{textblock*}{71.5mm}[0,1](144.5mm,-46.5mm)
    \textblockcolor{black} \rule{0mm}{\ruleheight}
  \end{textblock*}}

%%% Titre
\begin{adjustwidth*}{\gauche}{-15mm}
  \sffamily\fontseries{ub}\selectfont
  \raggedright
  \vspace*{2cm}
  \thetitle
\end{adjustwidth*}

%%% Affichage de la bande
\banderecto
\cleardoublepage

%%% Page de garde
\begin{adjustwidth*}{\gauche}{-15mm}
  \sffamily\fontseries{ub}\selectfont
  \raggedright
  \vspace*{2cm}
  \thetitle \\
  \bfseries
  \vspace*{3cm}
  \theauthor
  \vspace*{\fill}
  \thedate
\end{adjustwidth*}

\clearpage

%%% Page de notices
\begingroup
\calccentering{\unitlength}
\begin{adjustwidth*}{\unitlength}{-\unitlength}
  \small
  \setlength{\parindent}{0pt}
  \setlength{\parskip}{\baselineskip}

  {\textcopyright} 2013 François Pelletier \\


  \includegraphics[height=7mm,keepaspectratio=true]{by}\;%
  \includegraphics[height=7mm,keepaspectratio=true]{sa} \\
  Cette création est mise à disposition selon le contrat
  Paternité-Partage des conditions initiales à l'identique 2.5 Canada
  disponible en ligne
  \url{http://creativecommons.org/licenses/by-sa/2.5/ca/} ou par
  courrier postal à Creative Commons, 171 Second Street, Suite 300,
  San Francisco, California 94105, USA.

  \textbf{Historique de publication}
  \vspace{-\baselineskip}
  \begin{tabbing}
    Décembre 2013:\quad\= Première édition
  \end{tabbing}

  \textbf{Code source} \\
  Le code source {\LaTeX} de ce document est disponible à l'adresse
  \begin{quote}
    \url{https://github.com/franc00018/ACT-2010-Exercices}
  \end{quote}
  ou en communiquant directement avec les auteurs.

  \vspace{1cm}

  % ISBN \ISBN \\
  % Dépôt légal -- Bibliothèque et Archives nationales du Québec, 2014 \\
  % Dépôt légal -- Bibliothèque et Archives Canada, 2014
\end{adjustwidth*}
\endgroup

%%% Retour à la police normale
\renewcommand{\sfdefault}{phv}

\clearpage

%%% Local Variables:
%%% mode: latex
%%% TeX-master: "exercices_series_chrono"
%%% coding: utf-8-unix
%%% End:

\addtocounter{section}{2}
\section{Modèles de volatilité stochastique}

\subsection{Variance du processus ARCH(2)}

On identifie d'abord la moyenne conditionnelle de $y_t$:
\begin{align}
E_{t-1}[y_t] &= E_{t-1}[a_0 + a_1 y_{t-1} + \epsilon_t] \\
&= a_0 + a_1 y_{t-1}
\end{align}
La variance conditionnelle peut alors s'obtenir en utilisant la définition habituelle:
\begin{align*}
V_{t-1}[y_t | y_{t-1}, y_{t-2}, \ldots] &= E_{t-1}[y_t - E_{t-1}[y_t]]^{2} \\
&= E_{t-1}[(a_0 + a_1 y_{t-1} + \epsilon_t)-(a_0 + a_1 y_{t-1})]^{2} \\
&= E_{t-1}[\epsilon_t]^{2} \\
&= E_{t-1}[\alpha_0 + \alpha_1\epsilon_{t-1}^2 + \alpha_2\epsilon_{t-2}^2] \\
&= \alpha_0 + \alpha_1\epsilon_{t-1}^2 + \alpha_2\epsilon_{t-2}^2
\end{align*}

La variance inconditionnelle s'obtient en trouvant la solution particulière pour $y_t$:
\begin{align*}
•y_t &= a_0 + a_1 y_{t-1} + \epsilon_t \\
&= (1+a_1)a_0 + a_1^2 y_{t-2} + a_1 \epsilon_{t-1} + \epsilon_t \\
&= \cdots \\
&= (a+a_1+a_2+a_3+\ldots)a_0 + \epsilon_t + a_1\epsilon_{t-1} + a_2\epsilon_{t-2} + \ldots \\
&= \frac{a_0}{1-a_1} + \sum_{i=0}^{\infty} a_1^{i}\epsilon_{t-i}
\end{align*}

On évalue la variance de cette dernière expression:
\begin{align*}
•Var[y_t] &= Var[\sum_{i=0}^{\infty} a_1^{i}\epsilon_{t-i}] \\
&= \sum_{i=0}^{\infty} a_1^{2i} Var[\epsilon_{t-i}] \\
&= \frac{\sigma^2}{1-a_1^2}
\end{align*}

À partir de la définition, on a que:
\begin{align*}
E[\epsilon_t^2] &= \alpha_0 + \alpha_1 E_[\epsilon_{t-1}^2] + \alpha_2 E[\epsilon_{t-2}^2].
\end{align*}

Comme la variance inconditionnelle de $\epsilon_t$ est identique à celle de $\epsilon_{t-1}$ et $\epsilon_{t-2}$, on peut affirmer que:
\begin{align*}
E[\epsilon_t^2] &= \frac{\alpha_0}{1-\alpha_1-\alpha_2} \\
&= \sigma^2.
\end{align*}

On obtient donc que la variance inconditionnelle de $y_t$ est
\begin{align*}
•Var[y_t] &= \frac{\alpha_0}{(1-\alpha_1-\alpha_2)(1-a_1^2)}.
\end{align*}

\clearpage
\includegraphics[height=7mm,keepaspectratio=true]{by-sa}\\%
Cette création est mise à disposition selon le contrat
\href{http://creativecommons.org/licenses/by-sa/2.5/ca/deed.fr}{%
  Paternité-Partage à l'identique 2.5 Canada} de Creative Commons
disponible à l'adresse \\
http://creativecommons.org/licenses/by-sa/2.5/ca/deed.fr \\

En vertu de ce contrat, vous êtes libre de :

\begin{itemize}
\item \textbf{partager} --- reproduire, distribuer et communiquer
  l'{\oe}uvre;
\item \textbf{remixer} --- adapter l'{\oe}uvre;
\item utiliser cette {\oe}uvre à des fins commerciales.
\end{itemize}

Selon les conditions suivantes:\\

  \begin{tabularx}{\linewidth}{@{}lX@{}}
    \raisebox{-9mm}[0mm][13mm]{%
      \includegraphics[height=11mm,keepaspectratio=true]{by}} &
    \textbf{Attribution} --- Vous devez attribuer l'{\oe}uvre de la
    manière indiquée par l'auteur de l'{\oe}uvre ou le titulaire des
    droits (mais pas d'une manière qui suggérerait qu'ils vous
    soutiennent ou
    approuvent votre utilisation de l'{\oe}uvre). \\
    \raisebox{-9mm}{\includegraphics[height=11mm,keepaspectratio=true]{sa}}
    & \textbf{Partage à l'identique} --- Si vous modifiez, transformez
    ou adaptez cette {\oe}uvre, vous n'avez le droit de distribuer
    votre création que sous une licence identique ou similaire à
    celle-ci.
  \end{tabularx}

%%% Local Variables: 
%%% mode: latex
%%% TeX-master: t
%%% End: 

\end{document}
